% Part: incompleteness
% Chapter: representability-in-q
% Section: undecidability

\documentclass[../../../include/open-logic-section]{subfiles}

\begin{document}

\olfileid{inc}{req}{und}
\olsection{અનિર્ણેયતા}

સિદ્ધાંત~$\Th{T}$ને આપણે
\emph{અનિર્ણેય} કહીએ, જો એવી
કોઈ સંગણકીય પ્રક્રિયા ન હોય
જે $\Th{T}$ની ભાષાના કોઈપણ
વાક્ય~$!A$ માટે ``શું
$\Th{T}$માં~$!A$ સાબિત થાય
છે?'' તેનો સાચો જવાબ
અચૂક રીતે અને સાન્ત પગલાંમાં
આપે. તેથી $\Th{Q}$ નિર્ણેય
હોય ત્યારે અને માત્ર ત્યારે જ
અંકગણિતની ભાષાનું વાક્ય~$!A$
આપતાં, $\Th{Q} \Proves !A$
છે કે નહીં તે નક્કી કરતી
સંગણકીય પ્રક્રિયા હોય.
આને વધુ ચોક્કસ રીતે પૂછી
શકીએ: સંબંધ~$\Prov[\Th{Q}](y)$
$y$ માટે ત્યારે અને માત્ર ત્યારે જ
સાચો હોય કે $y$
$\Th{Q}$માં સાબિત થઈ શકે
એવા વાક્યની ગડેલ સંખ્યા
હોય, તો શું આ સંબંધ પુનરાવર્તી
છે? જવાબ છે: ના.

\begin{thm}
$\Th{Q}$ અનિર્ણેય છે, એટલે
આ સંબંધ
\[
\Prov[\Th{Q}](y) \defiff \fn{Sent}(y) \land
\lexists[x][\Prf[\Th{Q}](x, y)]
\]
પુનરાવર્તી નથી.
\end{thm}

\begin{proof}
ધારો કે તે પુનરાવર્તી હોય.
તો અટકવાની સમસ્યા આ રીતે
ઉકેલી શકીએ: $e$ અને $n$
આપેલાં હોય, તો આપણે જાણીએ
છીએ કે $\cfind{e}(n) \fdefined$
ત્યારે અને માત્ર ત્યારે જ
કોઈ~$s$ માટે $T(e, n, s)$ છે;
અહીં $T$ એ
\olref[cmp][rec][nft]{thm:kleene-nf}માંથી
ક્લીનીનું વિધેયધર્મ છે.
$T$ આદિમ પુનરાવર્તી હોવાથી
તેને સૂત્ર~$!B_T$ વડે
$\Th{Q}$માં નિરૂપી શકાય છે,
એટલે $T(e, n, s)$ હોય
ત્યારે અને માત્ર ત્યારે જ
$\Th{Q} \Proves
!B_T(\num{e}, \num{n}, \num{s})$.
જો $\Th{Q} \Proves
!B_T(\num{e}, \num{n}, \num{s})$,
તો $ \Th{Q} \Proves
\lexists[y][!B_T(\num{e}, \num{n}, y)]$
પણ છે. જો આવું કોઈ~$s$
ન હોય, તો દરેક~$s$ માટે
$\Th{Q} \Proves \lnot
!B_T(\num{e}, \num{n}, \num{s})$.
પરંતુ $\Th{Q}$
$\omega$-સુસંગત છે, એટલે
દરેક~$n \in \Nat$ માટે
$\Th{Q} \Proves \lnot
!A(\num{n})$ હોય, તો
$\Th{Q} \Proves/
\lexists[y][!A(y)]$.
આપણે આ જાણીએ છીએ, કારણ કે
$\Th{Q}$ની સ્વયંસિદ્ધિઓ
પ્રમાણભૂત નમૂના~$\Struct{N}$માં
સાચી છે. તેથી
$\Th{Q} \Proves/
\lexists[y][!B_T(\num{e}, \num{n}, y)]$.
બીજા શબ્દોમાં,
$\Th{Q} \Proves
\lexists[y][!B_T(\num{e}, \num{n}, y)]$
ત્યારે અને માત્ર ત્યારે જ
કોઈ~$s$ માટે $T(e, n, s)$,
એટલે ત્યારે અને માત્ર ત્યારે જ
$\cfind{e}(n) \fdefined$.
$e$ અને~$n$ પરથી આપણે
$\Gn{\lexists[y][!B_T(\num{e},
    \num{n}, y)]}$ ગણી શકીએ;
આ કામ કરતું આદિમ પુનરાવર્તી
વિધેય $g(e, n)$ લો. તો
\[
h(e, n) =
\begin{cases}
1 & \text{જો $\Prov[\Th{Q}](g(e, n))$}\\
0 & \text{અન્યથા}.
\end{cases}
\]
આથી $\Prov[\Th{Q}]$
પુનરાવર્તી હોય તો~$h$ પણ
પુનરાવર્તી હોવાનું નીકળે.
પરંતુ
\olref[cmp][rec][hlt]{thm:halting-problem}
પ્રમાણે~$h$ પુનરાવર્તી નથી;
તેથી $\Prov[\Th{Q}]$ પણ
પુનરાવર્તી ન હોઈ શકે.
\end{proof}

\begin{cor}
પ્રથમ-ક્રમનો તર્ક અનિર્ણેય છે.
\end{cor}

\begin{proof}
પ્રથમ-ક્રમનો તર્ક નિર્ણેય
હોય, તો $\Th{Q}$માં
સાબિતિક્ષમતા પણ નિર્ણેય
હોત, કારણ કે
$\Th{Q} \Proves !A$ ત્યારે અને
માત્ર ત્યારે જ
$\Proves !T \lif !A$;
અહીં~$!T$ એ $\Th{Q}$ની
સ્વયંસિદ્ધિઓનું સંયોજન છે.
\end{proof}

\end{document}
